
\documentclass{article}
\usepackage{colortbl}
\usepackage{makecell}
\usepackage{multirow}
\usepackage{supertabular}

\begin{document}

\newcounter{utterance}

\centering \large Interaction Transcript for game `wordle', experiment `high\_frequency\_words\_no\_clue\_no\_critic', episode 5 with mistral{-}large{-}3{-}azure.
\vspace{24pt}

{ \footnotesize  \setcounter{utterance}{1}
\setlength{\tabcolsep}{0pt}
\begin{supertabular}{c@{$\;$}|p{.15\linewidth}@{}p{.15\linewidth}p{.15\linewidth}p{.15\linewidth}p{.15\linewidth}p{.15\linewidth}}
    \# & \multicolumn{2}{c}{Player} && \multicolumn{2}{c}{Game Master} \\
    \hline

    \theutterance \stepcounter{utterance}  
    & & & \multicolumn{4}{p{0.6\linewidth}}{
        \cellcolor[rgb]{0.9,0.9,0.9}{
            \makecell[{{p{\linewidth}}}]{
                \texttt{\tiny{[P1$\langle$GM]}}
                \texttt{Vous êtes un magicien du langage qui aime deviner des mots en utilisant les règles données.} \\
\\ 
\texttt{Bienvenue à Wordle ! Vous avez six tentatives pour deviner le mot cible, un mot valide de cinq lettres minuscules. Veuillez utiliser les balises "explication :" et "proposition :" pour fournir une explication concise pour chaque proposition.} \\
\\ 
\texttt{Par exemple, si votre proposition est "année", votre réponse devrait être} \\
\texttt{explication : c'est un mot commun de cinq lettres, et je commence ma proposition avec ce mot.} \\
\texttt{proposition : année} \\
\\ 
\texttt{Après chaque proposition, votre réponse sera validée, et vous recevrez un retour indiquant quelles lettres sont correctes (vert), quelles lettres sont correctes mais à la mauvaise position (jaune), et quelles lettres sont incorrectes (rouge). Ce retour peut être utile pour déterminer quelles lettres inclure ou exclure dans votre prochaine proposition.} \\
\\ 
\texttt{Par exemple, le retour pour "année" pourrait être :} \\
\texttt{retour\_proposition : a<jaune> n<jaune> n<vert> é<jaune> e<rouge>} \\
\\ 
\texttt{L'explication doit contenir des détails sur la manière dont le retour\_proposition est utilisé pour arriver à une nouvelle proposition.} \\
\\ 
\texttt{Ne générez aucun autre texte.} \\
\\ 
\texttt{Commençons avec votre première proposition.} \\
            }
        }
    }
     \\ \\

    \theutterance \stepcounter{utterance}  
    & \multicolumn{4}{p{0.6\linewidth}}{
        \cellcolor[rgb]{0.561,0.737,0.561}{
            \makecell[{{p{\linewidth}}}]{
                \texttt{\tiny{[P1$\rangle$GM]}}
                \texttt{explication : Je commence par un mot courant avec des voyelles fréquentes et des consonnes variées pour maximiser l'information sur les lettres présentes.} \\
\texttt{proposition : audio} \\
            }
        }
    }
    & & \\ \\

    \theutterance \stepcounter{utterance}  
    & & & \multicolumn{2}{p{0.3\linewidth}}{
        \cellcolor[rgb]{0.9,0.9,0.9}{
            \makecell[{{p{\linewidth}}}]{
                \texttt{\tiny{[GM$|$GM]}}
                \texttt{attempts = 1} \\
\texttt{target\_word = jeudi} \\
\texttt{guess = audio} \\
\texttt{guess\_feedback = a<red> u<yellow> d<yellow> i<yellow> o<red>} \\
            }
        }
    }
    & & \\ \\

    \theutterance \stepcounter{utterance}  
    & & & \multicolumn{4}{p{0.6\linewidth}}{
        \cellcolor[rgb]{0.9,0.9,0.9}{
            \makecell[{{p{\linewidth}}}]{
                \texttt{\tiny{[P1$\langle$GM]}}
                \texttt{retour\_proposition : a<rouge> u<jaune> d<jaune> i<jaune> o<rouge>} \\
\\ 
\texttt{Fournissez votre réponse uniquement dans ce format.} \\
\texttt{explication : détails} \\
\texttt{proposition : mot} \\
            }
        }
    }
     \\ \\

    \theutterance \stepcounter{utterance}  
    & \multicolumn{4}{p{0.6\linewidth}}{
        \cellcolor[rgb]{0.561,0.737,0.561}{
            \makecell[{{p{\linewidth}}}]{
                \texttt{\tiny{[P1$\rangle$GM]}}
                \texttt{explication : "a" et "o" sont absents (rouge), donc je les exclus. "u", "d" et "i" sont présents mais mal placés (jaune). Je teste un mot avec ces lettres dans d'autres positions, en évitant "a" et "o", et en ajoutant de nouvelles consonnes pour élargir les possibilités.} \\
\texttt{proposition : durit} \\
            }
        }
    }
    & & \\ \\

    \theutterance \stepcounter{utterance}  
    & & & \multicolumn{2}{p{0.3\linewidth}}{
        \cellcolor[rgb]{0.9,0.9,0.9}{
            \makecell[{{p{\linewidth}}}]{
                \texttt{\tiny{[GM$|$GM]}}
                \texttt{attempts = 2} \\
\texttt{target\_word = jeudi} \\
\texttt{guess = durit} \\
\texttt{guess\_feedback = d<yellow> u<yellow> r<red> i<yellow> t<red>} \\
            }
        }
    }
    & & \\ \\

    \theutterance \stepcounter{utterance}  
    & & & \multicolumn{4}{p{0.6\linewidth}}{
        \cellcolor[rgb]{0.9,0.9,0.9}{
            \makecell[{{p{\linewidth}}}]{
                \texttt{\tiny{[P1$\langle$GM]}}
                \texttt{retour\_proposition : d<jaune> u<jaune> r<rouge> i<jaune> t<rouge>} \\
\\ 
\texttt{Fournissez votre réponse uniquement dans ce format.} \\
\texttt{explication : détails} \\
\texttt{proposition : mot} \\
            }
        }
    }
     \\ \\

    \theutterance \stepcounter{utterance}  
    & \multicolumn{4}{p{0.6\linewidth}}{
        \cellcolor[rgb]{0.561,0.737,0.561}{
            \makecell[{{p{\linewidth}}}]{
                \texttt{\tiny{[P1$\rangle$GM]}}
                \texttt{explication : "r" et "t" sont absents (rouge). "d", "u" et "i" sont présents mais mal placés (jaune). Je sais que :} \\
\texttt{{-} "d" n'est pas en 1ère position, "u" pas en 2ème, "i" pas en 4ème.} \\
\texttt{Je propose un mot avec ces lettres dans d'autres positions, en évitant "r", "t", "a", "o", et en testant une nouvelle consonne ("l").} \\
\texttt{proposition : lundi} \\
            }
        }
    }
    & & \\ \\

    \theutterance \stepcounter{utterance}  
    & & & \multicolumn{2}{p{0.3\linewidth}}{
        \cellcolor[rgb]{0.9,0.9,0.9}{
            \makecell[{{p{\linewidth}}}]{
                \texttt{\tiny{[GM$|$GM]}}
                \texttt{attempts = 3} \\
\texttt{target\_word = jeudi} \\
\texttt{guess = lundi} \\
\texttt{guess\_feedback = l<red> u<yellow> n<red> d<green> i<green>} \\
            }
        }
    }
    & & \\ \\

    \theutterance \stepcounter{utterance}  
    & & & \multicolumn{4}{p{0.6\linewidth}}{
        \cellcolor[rgb]{0.9,0.9,0.9}{
            \makecell[{{p{\linewidth}}}]{
                \texttt{\tiny{[P1$\langle$GM]}}
                \texttt{retour\_proposition : l<rouge> u<jaune> n<rouge> d<vert> i<vert>} \\
\\ 
\texttt{Fournissez votre réponse uniquement dans ce format.} \\
\texttt{explication : détails} \\
\texttt{proposition : mot} \\
            }
        }
    }
     \\ \\

    \theutterance \stepcounter{utterance}  
    & \multicolumn{4}{p{0.6\linewidth}}{
        \cellcolor[rgb]{0.561,0.737,0.561}{
            \makecell[{{p{\linewidth}}}]{
                \texttt{\tiny{[P1$\rangle$GM]}}
                \texttt{explication :} \\
\texttt{{-} "l" et "n" sont absents (rouge), je les exclus.} \\
\texttt{{-} "d" est bien placé en 4ème position (vert) et "i" en 5ème (vert).} \\
\texttt{{-} "u" est présent mais mal placé (jaune), donc pas en 2ème position.} \\
\texttt{Je cherche un mot avec "d" en 4ème, "i" en 5ème, "u" ailleurs, et sans "a", "o", "r", "t", "l", "n".} \\
\texttt{proposition : jouir} \\
            }
        }
    }
    & & \\ \\

    \theutterance \stepcounter{utterance}  
    & & & \multicolumn{2}{p{0.3\linewidth}}{
        \cellcolor[rgb]{0.9,0.9,0.9}{
            \makecell[{{p{\linewidth}}}]{
                \texttt{\tiny{[GM$|$GM]}}
                \texttt{attempts = 4} \\
\texttt{target\_word = jeudi} \\
\texttt{guess = jouir} \\
\texttt{guess\_feedback = j<green> o<red> u<green> i<yellow> r<red>} \\
            }
        }
    }
    & & \\ \\

    \theutterance \stepcounter{utterance}  
    & & & \multicolumn{4}{p{0.6\linewidth}}{
        \cellcolor[rgb]{0.9,0.9,0.9}{
            \makecell[{{p{\linewidth}}}]{
                \texttt{\tiny{[P1$\langle$GM]}}
                \texttt{retour\_proposition : j<vert> o<rouge> u<vert> i<jaune> r<rouge>} \\
\\ 
\texttt{Fournissez votre réponse uniquement dans ce format.} \\
\texttt{explication : détails} \\
\texttt{proposition : mot} \\
            }
        }
    }
     \\ \\

    \theutterance \stepcounter{utterance}  
    & \multicolumn{4}{p{0.6\linewidth}}{
        \cellcolor[rgb]{0.561,0.737,0.561}{
            \makecell[{{p{\linewidth}}}]{
                \texttt{\tiny{[P1$\rangle$GM]}}
                \texttt{explication :} \\
\texttt{{-} "j" est bien placé en 1ère position (vert) et "u" en 3ème (vert).} \\
\texttt{{-} "i" est présent mais mal placé (jaune), donc pas en 5ème (déjà confirmé comme vert dans la proposition précédente → contradiction, donc "i" doit être ailleurs).} \\
\texttt{{-} "o" et "r" sont absents (rouge).} \\
\texttt{Je cherche un mot avec :} \\
\texttt{  {-} "j" en 1, "u" en 3, "d" en 4, "i" en 2 ou 5 (mais 5 était vert avant → donc en 2).} \\
\texttt{  {-} Pas de "a", "o", "r", "t", "l", "n".} \\
\texttt{proposition : judas} \\
            }
        }
    }
    & & \\ \\

    \theutterance \stepcounter{utterance}  
    & & & \multicolumn{2}{p{0.3\linewidth}}{
        \cellcolor[rgb]{0.9,0.9,0.9}{
            \makecell[{{p{\linewidth}}}]{
                \texttt{\tiny{[GM$|$GM]}}
                \texttt{attempts = 5} \\
\texttt{target\_word = jeudi} \\
\texttt{guess = judas} \\
\texttt{guess\_feedback = j<green> u<yellow> d<yellow> a<red> s<red>} \\
            }
        }
    }
    & & \\ \\

    \theutterance \stepcounter{utterance}  
    & & & \multicolumn{4}{p{0.6\linewidth}}{
        \cellcolor[rgb]{0.9,0.9,0.9}{
            \makecell[{{p{\linewidth}}}]{
                \texttt{\tiny{[P1$\langle$GM]}}
                \texttt{retour\_proposition : j<vert> u<jaune> d<jaune> a<rouge> s<rouge>} \\
\\ 
\texttt{Fournissez votre réponse uniquement dans ce format.} \\
\texttt{explication : détails} \\
\texttt{proposition : mot} \\
            }
        }
    }
     \\ \\

    \theutterance \stepcounter{utterance}  
    & \multicolumn{4}{p{0.6\linewidth}}{
        \cellcolor[rgb]{0.561,0.737,0.561}{
            \makecell[{{p{\linewidth}}}]{
                \texttt{\tiny{[P1$\rangle$GM]}}
                \texttt{explication :} \\
\texttt{{-} "j" est bien placé en 1ère position (vert).} \\
\texttt{{-} "u" est présent mais pas en 2ème (jaune), donc doit être en 3ème (déjà confirmé vert précédemment → contradiction, donc erreur dans l'analyse précédente).} \\
\texttt{{-} "d" est présent mais pas en 3ème (jaune), donc doit être en 4ème (vert dans "lundi").} \\
\texttt{{-} "a" et "s" sont absents (rouge).} \\
\texttt{{-} "i" doit être en 5ème (vert dans "lundi"), mais "jouir" montrait "i" mal placé → incohérence. Je réévalue :} \\
\texttt{  {-} Mot cible : "j" en 1, "u" en 3, "d" en 4, "i" en 5 (mais "jouir" montrait "i" en 5 jaune → donc "i" n'est pas en 5).} \\
\texttt{  {-} Nouvelle hypothèse : "i" en 2, "u" en 3, "d" en 4, et lettre inconnue en 5 (sans "a", "o", "r", "t", "l", "n", "s").} \\
\texttt{proposition : jupes} \\
            }
        }
    }
    & & \\ \\

    \theutterance \stepcounter{utterance}  
    & & & \multicolumn{2}{p{0.3\linewidth}}{
        \cellcolor[rgb]{0.9,0.9,0.9}{
            \makecell[{{p{\linewidth}}}]{
                \texttt{\tiny{[GM$|$GM]}}
                \texttt{Error: The guessed word is not a valid word for this game.} \\
            }
        }
    }
    & & \\ \\

    \theutterance \stepcounter{utterance}  
    & & & \multicolumn{4}{p{0.6\linewidth}}{
        \cellcolor[rgb]{0.9,0.9,0.9}{
            \makecell[{{p{\linewidth}}}]{
                \texttt{\tiny{[P1$\langle$GM]}}
                \texttt{Votre proposition n'est pas un mot valide pour ce jeu. Veuillez réessayer.} \\
\\ 
\texttt{Fournissez votre réponse uniquement dans ce format.} \\
\texttt{explication : détails} \\
\texttt{proposition : mot} \\
            }
        }
    }
     \\ \\

    \theutterance \stepcounter{utterance}  
    & \multicolumn{4}{p{0.6\linewidth}}{
        \cellcolor[rgb]{0.561,0.737,0.561}{
            \makecell[{{p{\linewidth}}}]{
                \texttt{\tiny{[P1$\rangle$GM]}}
                \texttt{explication :} \\
\texttt{{-} Le mot doit contenir "j" en 1ère position (vert), "d" en 4ème (vert), et "i" quelque part (jaune dans "lundi" et "jouir").} \\
\texttt{{-} "u" est présent mais pas en 2ème (jaune dans "judas"), donc en 3ème.} \\
\texttt{{-} Lettres exclues : "a", "o", "r", "t", "l", "n", "s".} \\
\texttt{{-} Je teste un mot avec "j" en 1, "u" en 3, "d" en 4, et "i" en 2 ou 5.} \\
\texttt{proposition : jadis} \\
            }
        }
    }
    & & \\ \\

    \theutterance \stepcounter{utterance}  
    & & & \multicolumn{2}{p{0.3\linewidth}}{
        \cellcolor[rgb]{0.9,0.9,0.9}{
            \makecell[{{p{\linewidth}}}]{
                \texttt{\tiny{[GM$|$GM]}}
                \texttt{attempts = 6} \\
\texttt{target\_word = jeudi} \\
\texttt{guess = jadis} \\
\texttt{guess\_feedback = j<green> a<red> d<yellow> i<yellow> s<red>} \\
            }
        }
    }
    & & \\ \\

    \theutterance \stepcounter{utterance}  
    & & & \multicolumn{2}{p{0.3\linewidth}}{
        \cellcolor[rgb]{0.9,0.9,0.9}{
            \makecell[{{p{\linewidth}}}]{
                \texttt{\tiny{[GM$|$GM]}}
                \texttt{game\_result = LOSS} \\
            }
        }
    }
    & & \\ \\

\end{supertabular}
}

\end{document}
